\chapter{Introduction}\label{chapter:introduction}

\section{Context and Motivation}

Robots operate in changing environments. Their behaviour depends on software,
hardware, sensors, communication, and the physical world. They also increasingly
use artificial-intelligence components for perception and decision making.
Uncertain sensor input, model confidence, and conditions that differ from
training data can therefore influence a robot after deployment. Testing remains
essential, but it cannot cover every environment and interaction that a robot
may encounter. Developers also need to observe the robot during operation.

Many modern robots are built with the Robot Operating System~2 (ROS~2).
Despite its name, ROS~2 is not an operating system in the traditional sense;
it runs on top of a general-purpose operating system such as Linux. It
provides software libraries and tools for assembling the functional components
of a robotic system, and reusable packages that work together in different
robot applications~\cite{macenski2022ros2}. A ROS~2 application is composed of
\emph{nodes}: programs with a focused task, such as reading a sensor or
planning a route. Nodes communicate through typed interfaces. \emph{Topics}
are mainly used for continuous data streams, \emph{services} for
request--response interactions, and \emph{actions} for long-running tasks that
provide feedback while they run~\cite{macenski2022ros2,ros2interfaces}.
Chapter~\ref{chapter:background} explains these mechanisms in detail.

A common way to observe a system during operation is \emph{runtime
monitoring}. Runtime monitoring collects information from a running system;
the collected information can be used to find errors, understand runtime
behaviour, or check whether the system follows expected
properties~\cite{rabiser2019domain,leucker2009brief}. A monitor may report,
for example, that commanded velocity exceeded a limit, a service call did not
have the expected effect, or an action missed a deadline. Unlike model
checking, runtime monitoring does not explore every possible execution path;
it analyses behaviour that is observed in actual
executions~\cite{leucker2009brief}.

Runtime monitoring is particularly relevant for ROS~2 applications. The
modular node structure makes it easier to reuse and replace software parts,
but it also spreads the behaviour of the robot across many nodes and
interfaces. Important information may come from several topics, services, or
actions, so it can be difficult to understand what the full system is doing by
looking at any single component. At the same time, ROS~2 offers several useful
observation points: a monitor can subscribe to a topic, read service
introspection events, or observe the hidden topics used by actions, and lower
layers expose network and execution traces~\cite{macenski2022ros2}. The
challenge is to turn these observation mechanisms into a reusable and
configurable monitoring infrastructure. The next section states this problem
more precisely.

\section{Problem Statement}\label{sec:problem}

A monitoring solution must make several connected choices. The developer has
to decide \emph{what} to observe: which topics, services, and actions carry
the relevant behaviour. They have to decide \emph{how} to observe it: through
extra subscriptions, service introspection, lower-level tracing, or
instrumented application code. They have to decide how much data to keep:
complete messages support offline replay, while a few fields from a high-rate
stream may be all that a check needs. They have to decide how the properties
of interest are specified and checked: as simple thresholds, state machines,
temporal-logic formulas, or project-specific checking services. Finally, they
have to decide \emph{where} the monitoring components are deployed: together
with the application, close to the robot, or on a separate machine that
receives the observations over a network.

These choices have advantages and disadvantages, and they interact with each
other. Observing everything makes checking flexible but can overload storage
and transport; a remote verifier eases the load on the robot but adds a
network boundary. The choices also change between applications and
deployments: a local developer test may keep collection and checking in one
process, while a fielded robot may keep ROS~2-facing collection near the robot
and send records to a separate verifier. Because the choices interact, an
ad-hoc combination assembled for one application rarely carries over to the
next one. There is a need for a systematic way to specify a monitoring
solution: a description of the available options and a means to select among
them.

The second problem is that implementing a chosen monitoring solution is time
consuming. Existing approaches address important parts of this problem, but
they use different observation levels, data formats, property languages, and
deployment assumptions. Much of the required code---subscribing to interfaces,
decoding typed messages, attaching timestamps and context, serializing data,
and transporting it to checkers and files---is essentially the same in every
project, yet it is written again for each application. This boilerplate
obscures the application-specific decisions and makes monitoring setups harder
to reproduce. The repeated implementation effort should be reduced by
automation and by reusable, configurable infrastructure support.

These two problems, the systematic specification of a monitoring solution and
its reusable implementation, are the subject of this thesis. Property
specification and property checking are themselves important research
problems: temporal logics, state machines, grammars, and rule languages
provide different solutions, and the runtime-verification literature surveys
this design space~\cite{leucker2009brief,francalanza2018distributed}. This
thesis does not propose a new property language or checking algorithm.
Instead, it focuses on the infrastructure that supplies observations to a
selected checker. A checker-specific converter and verdict service are treated
as extension points.

The problem addressed by this thesis is therefore:

\begin{quote}
How can heterogeneous ROS~2 observations be collected, represented, processed,
routed, and delivered to property checkers through a reusable infrastructure
whose deployment and checker-specific parts remain configurable?
\end{quote}

\section{Objective and Research Questions}

The objective is to design, implement, and evaluate a configurable monitoring
infrastructure for ROS~2 applications. The infrastructure should cover typed
topic messages, service events, and action feedback and status; preserve enough
context to relate a verdict to its observations; and support both integrated
and split placement of collection and checking.

The objective is refined into four research questions:

\begin{description}
  \item[RQ1.] What monitoring capabilities and ROS~2 observation mechanisms
  exist, and which application-level observation gaps remain in the reviewed
  approaches?

  \item[RQ2.] What are the commonalities and variabilities of monitoring
  solutions for ROS~2 applications?

  \item[RQ3.] How can a reference architecture incorporate the selected
  monitoring features while separating reusable infrastructure from
  property-specific checking?

  \item[RQ4.] To what extent do the proposed reference architecture and its
  implementation support the monitoring needs of representative ROS~2
  applications?
\end{description}

\section{Scope}

The prototype follows a deliberately focused path through the monitoring design
space. It observes the application-level data that ROS~2 exposes through
topics, service introspection, and hidden action feedback and status topics.
It supports passive online monitoring and recorded-data replay. Checking may run
inside the ROS~2-facing monitor or in a separate runtime reached through a file
or MQTT transport.

The following concerns are outside the implemented scope: a new property
formalism, automatic synthesis of checking algorithms, active intervention in
the system under monitoring, global ordering across hosts, fault-tolerant
distributed verification, network-packet inspection, and operating-system
tracing. The reference architecture keeps explicit extension points for several
of these concerns; Chapter~\ref{chapter:discussion} discusses which extension
points remain open.

\section{Approach and Contributions}

Each research question is answered as follows.

To answer RQ1, existing research and tools for monitoring ROS and ROS~2
applications are reviewed using consistent criteria. The approaches are
grouped by the level at which they observe the system: the ROS~2 application
interfaces, network-level observation, and lower-level tracing. The review
identifies which ROS~2 elements each approach can observe and which
observation gaps remain (Chapters~\ref{chapter:background}
and~\ref{chapter:related_work}).

To answer RQ2, a domain analysis of the reviewed monitoring approaches is
performed. It identifies the tasks they share, such as collecting, processing,
routing, checking, and reporting, and the points where they differ, such as
the observable elements, the monitoring architecture, the monitoring mode, the
system configuration, and the implementation paradigm. These commonalities and
variabilities are captured in a feature model (Chapter~\ref{chapter:design}).

To answer RQ3, a reference architecture is derived from the feature model and
implemented as a Python prototype (Chapters~\ref{chapter:design}
and~\ref{chapter:implementation}). The architecture is a configurable
pipeline: collectors observe configured topics, service event streams, or
action feedback and status streams and store each observation in a common data
record; transformers reduce or reshape records; dispatchers route them to
outputs and checking components; converters prepare them for property-specific
verdict services; and exporters deliver data and verdicts to files, standard
output, or MQTT. The same pipeline supports an integrated deployment, in which
all steps run in one ROS~2-facing process, and a split deployment, in which a
robot-side monitor sends records over MQTT to a separate verifier without a
ROS~2 dependency.

To answer RQ4, five experiments evaluate the selected feature configuration,
progressing from a controlled topic property to a Navigation~2
(Nav2)~\cite{macenski2020marathon} case
study, a two-robot correlation property, and a deployment across three
computers. Nav2 is chosen because it combines topics, services, and actions in
an existing, widely used robotic application. The evaluation also checks
whether the reference architecture can accommodate the principal capabilities
of reviewed monitoring approaches (Chapter~\ref{chapter:evaluation}).

The contributions are:

\begin{enumerate}
  \item a domain analysis and feature model for ROS~2 monitoring
  solutions;
  \item a reference architecture that separates observation, common data
  representation, processing, routing, checker adaptation, checking, and
  result use;
  \item a prototype that collects topic messages, service request and response
  events, and action feedback and status into one common record format;
  \item plug-in interfaces for checker-specific converters and verdict
  services, with integrated and split deployment;
  \item a deployment-specification generator that creates matching runtime YAML
  files and reduces low-level, routine monitoring boilerplate; and
  \item an evaluation using controlled stimuli and Nav2 case studies, including
  evidence that the architecture can accommodate the main capabilities of four
  reviewed monitoring approaches.
\end{enumerate}

The main contribution is the reusable monitoring architecture and the
infrastructure that realizes it. The thesis provides the architectural layer
between ROS~2 communication and property-specific checking. This layer can
host deployment styles used in existing monitoring frameworks and extend them
with uniform ROS~2 observations, split robot/verifier deployment, replay, and
evidence handling.

\section{Thesis Structure}

Chapter~\ref{chapter:background} introduces ROS~2 communication, monitoring
concepts, observation mechanisms, feature models, and deployment terminology.
Chapter~\ref{chapter:related_work} reviews the selected approaches and derives
the research gaps. Chapter~\ref{chapter:design} presents the feature model and
reference architecture. Chapter~\ref{chapter:implementation} explains the
prototype and configuration workflow. Chapter~\ref{chapter:evaluation}
evaluates the implementation and the architectural accommodation of related
work. Chapter~\ref{chapter:discussion} discusses the research questions,
trade-offs, limitations, threats to validity, and future work.
Chapter~\ref{chapter:conclusions} gives concise final answers and conclusions.
